\section{实验评估}
\label{sec:evaluation}

本章按照端到端执行、执行前控制和可复现失败记录三个原则评估AgentGate。与将已生成轨迹逐条
输入分类器的离线评估不同，本文的主实验从任务开始运行Agent或确定性任务执行器，实际调用隔离
工具，并由环境中的副作用判定器检查消息、文件、命令和安装操作是否真实发生。只有在危险副作用
发生前终止调用才计为防护成功。实验回答以下五个研究问题：

\hangindent=2.5em\hangafter=1\textbf{RQ1（端到端有效性）}：AgentGate能否阻止实际攻击任务，
同时保持正常任务可用？

\hangindent=2.5em\hangafter=1\textbf{RQ2（状态化机制）}：事件序列、ATG数据来源关系和安全标签
传播分别对多步风险检测有何贡献？

\hangindent=2.5em\hangafter=1\textbf{RQ3（公开基线）}：在相同Agent、任务环境和攻击设置下，
AgentGate与可执行的公开运行时防护工具相比表现如何？

\hangindent=2.5em\hangafter=1\textbf{RQ4（开销与扩展性）}：运行时开销如何随工具调用数量、
ATG节点、边和来源深度增长？

\hangindent=2.5em\hangafter=1\textbf{RQ5（模型鲁棒性）}：固定任务Agent，仅替换语义解析模型时，
最终决策、任务效用、延迟和Token消耗是否稳定？

\subsection{实验设置}
\label{subsec:evaluation-setup}

\textbf{运行环境。}
本轮实验对应AgentGate 0.6.0和Python 3.12.3，运行于Linux 6.8服务器。服务器包含两颗
AMD EPYC 9554处理器，共128个物理核心和503\,GiB内存。默认任务模型和默认语义模型均为
DeepSeek-V4-Pro-0813；RQ5固定确定性任务Agent，只改变语义模型。API密钥和服务URL只从
\texttt{.env}读取，不写入实验结果。所有任务级和调用级结果以统一JSONL保存，并记录benchmark
revision、AgentGate commit、配置摘要、模型、调用数、执行决策、真实副作用、延迟、Token和ATG
快照。

\textbf{Benchmark。}
本文实际执行四组工作负载。第一，AgentGate-StatefulBench包含8个攻击及8个序列相似的良性
配对任务，覆盖敏感读取后外发、编码转换后外发、文件中转、不可信内容执行与安装、不可信内容
驱动的敏感读取、渐进式收集以及跨MCP Server传播。每个任务运行真实工具实现，并使用独立临时
目录保存副作用。第二，AgentDojo固定为commit \texttt{089ed468}、版本0.1.35和benchmark
1.2，本轮运行\texttt{workspace}中的\texttt{user\_task\_0--1}与
\texttt{tool\_knowledge/injection\_task\_0}组合。第三，规模工作负载实际执行5、10、20、40和
80次工具调用，逐步增加数据对象和来源深度。第四，语义端到端工作负载包含一个模糊命名的
\texttt{relay\_record}工具；攻击任务将真实敏感读取结果交给外部relay，配对正常任务则传递
公开数据，每个模型重复3次。

AgentDojo只执行了2个攻击组合，因此仅用于验证真实框架接入和任务兼容性，不足以估计总体ASR。
MSB固定commit \texttt{c1d6a701}，其本地filesystem MCP Server已成功启动，但Paper Search
Server跳转至交互式Smithery OAuth并等待浏览器回调，未产生可计分任务。MCP-SafetyBench会对
真实GitHub仓库、文件系统、浏览器和终端产生操作；当前环境没有完整的专用一次性账户与隔离凭证，
因此未执行。MCP-Bench缺少其配置要求的模型、LLM judge及MCP Server密钥，也未执行。上述三项
均以\texttt{not\_run}而非零分记录，且不进入主结果分母。

\textbf{Baseline和消融。}
AgentDojo比较No Defense、官方Tool Filter和AgentGate。由于官方OpenAI-compatible路径将
provider名称错误用作模型ID，适配器仅将filter请求的model字段替换为固定任务模型ID，不改变
过滤策略。AgentGate适配器替换AgentDojo的\texttt{ToolsExecutor}，原环境和官方utility/security
evaluator保持不变。AgentGuard 0.9.0（commit \texttt{3d235041}）、Invariant Guardrails
（\texttt{2340fe2d}）和MCPKernel 0.3.0（\texttt{bdab6f57}）的源码与版本均已固定，但由于
可执行MCP工作负载在认证预检阶段失败，本文不对它们构造离线替代分数，也不修改其核心设计以填满
比较矩阵。

RQ2采用五级消融：A0只检查当前ToolSecurityEvent；A1保留事件顺序但不建立来源关系；A2保留
Agent、ToolEvent、Resource、TrustDomain及结构边，但删除DataObject和
\texttt{PRODUCES/CONSUMES/DERIVES\_FROM}；A3建立provenance但关闭派生对象的标签传播；A4为
完整AgentGate。为隔离所研究机制，RQ2去掉默认外部SEND审批等通用兜底策略，各模式运行完全相同
的16个任务。该设置是机制消融，不等同于生产策略配置。

\textbf{指标。}
攻击成功率ASR是实际产生有害副作用的攻击占比；防护成功率为$1-ASR$；良性完成率BCR是未被
错误阻断且完成预期任务的良性任务占比。另报告真实副作用率、晚检测率和配对误报率。RQ1比例给出
Wilson 95\%置信区间。RQ4使用\texttt{perf\_counter}记录每次调用的墙钟时间，报告P50、P95、
P99、吞吐和图的序列化内存估计。RQ5将HTTP成功、严格Schema成功和最终执行决策分开；Schema
失败执行失败关闭，因此会降低BCR并计入良性拒绝，同时单独报告失败率。

\subsection{RQ1：端到端安全性与正常任务效用}
\label{subsec:rq1-e2e}

表~\ref{tab:rq1-e2e}给出实际执行结果。StatefulBench中No Defense的8个攻击全部产生了有害
副作用；完整AgentGate在sink执行前阻断8/8个攻击，未出现晚检测。其8个良性配对任务完成7个，
唯一失败为“不可信内容读取后，读取公开数据并外发”的负对照。系统依据
\texttt{untrusted\_context\_high\_risk}要求审批，说明当前控制关联规则在缺少显式数据依赖时
仍较保守。

\begin{table}[htbp]
    \centering
    \caption{端到端攻击成功率与正常任务完成率}
    \label{tab:rq1-e2e}
    \resizebox{0.94\linewidth}{!}{
    \begin{tabular}{ll|rrr|rr}
        \toprule
        Benchmark & Defense & 攻击数 & ASR(\%) & 副作用数 & 良性数 & BCR(\%) \\
        \midrule
        StatefulBench & No Defense & 8 & 100.0 & 8 & 8 & 100.0 \\
        StatefulBench & Full AgentGate & 8 & 0.0 & 0 & 8 & 87.5 \\
        AgentDojo subset & No Defense & 2 & 0.0 & 0 & 2 & 100.0 \\
        AgentDojo subset & Tool Filter & 2 & 0.0 & 0 & 2 & 0.0 \\
        AgentDojo subset & AgentGate & 2 & 0.0 & 0 & 2 & 100.0 \\
        \bottomrule
    \end{tabular}}
\end{table}

StatefulBench的ASR从8/8降至0/8，其Wilson 95\%区间分别为[67.6\%,100\%]和
[0\%,32.4\%]；完整系统BCR为7/8，区间为[52.9\%,97.8\%]。区间较宽表明该集合用于验证
机制，而非提供精确的开放世界总体估计。

AgentDojo的三个配置均未完成注入目标，包括No Defense，因此不能将ASR为零归因于防护。
Tool Filter在两个任务中均移除了正常任务所需工具，BCR为零；AgentGate和No Defense均完成
两个正常任务。AgentGate在组合任务中没有触发阻断，说明本轮结果只证明其能位于真实AgentDojo
执行边界且不破坏这两个任务，不能证明相对安全增益。后续必须扩大user task、injection task、
suite和随机种子覆盖。

\textbf{RQ1结论。}
完整AgentGate在受控的8类真实副作用链上实现执行前阻断，代价是1/8的配对良性任务需要审批。
公开AgentDojo小样本验证了接入兼容性，但因No Defense本身未受攻击而不构成有效安全对照。

\subsection{RQ2：ATG、来源关系与标签传播}
\label{subsec:rq2-stateful}

表~\ref{tab:rq2-ablation}显示，A0无法识别任何多步攻击。A1仅依据“曾读取敏感数据后又发送”或
“曾读取不可信内容后又执行”的先后顺序，可阻断全部攻击，但同时误阻断5/8个序列相似的良性任务。
这说明简单时序共现具有高召回但缺少数据一致性证据。

\begin{table}[htbp]
    \centering
    \caption{StatefulBench状态化机制消融}
    \label{tab:rq2-ablation}
    \resizebox{0.96\linewidth}{!}{
    \begin{tabular}{l|rrrr|rrr}
        \toprule
        配置 & ASR(\%) & BCR(\%) & 配对FPR(\%) & 攻击阻断 & Prov.阻断 & 来源边 & 最大深度 \\
        \midrule
        A0 Event-only & 100 & 100.0 & 0.0 & 0/8 & 0/8 & 0 & 0 \\
        A1 Event + Sequence & 0 & 37.5 & 62.5 & 8/8 & 8/8 & 0 & 0 \\
        A2 ATG w/o Provenance & 50 & 87.5 & 12.5 & 4/8 & 4/8 & 0 & 0 \\
        A3 Provenance w/o Labels & 25 & 87.5 & 12.5 & 6/8 & 6/8 & 11 & 1 \\
        A4 Full AgentGate & 0 & 87.5 & 12.5 & 8/8 & 8/8 & 9 & 1 \\
        \bottomrule
    \end{tabular}}
\end{table}

A2仍可利用资源、目标、时间窗口和不可信上下文等非provenance结构阻断4类攻击，但漏掉直接外发、
转换外发、文件中转和跨Server传播。A3恢复直接值和简单转换的数据依赖后阻断6类攻击，但文件中转
与跨Server转换需要标签沿派生对象持续传播，因而仍漏报。A4最终阻断8类攻击。分模式实际构造的
来源边数量不能直接横向解释为覆盖率，因为更早的阻断会减少后续事件提交；A4的阻断率和配对误报
才是主要终点指标。

进一步按风险链检查，A4对Sensitive Read$\rightarrow$Send、Read$\rightarrow$Transform
$\rightarrow$Send、Read$\rightarrow$Write$\rightarrow$Send、Untrusted Read$\rightarrow$
Execute/Install、外部控制$\rightarrow$敏感读取$\rightarrow$外发、累计读取$\rightarrow$批量
外发和跨Server转换外发均为1/1阻断。A1也全部命中，但其62.5\%配对误报证明了“操作出现过”
不能替代“当前sink实际消费了该数据”。

\textbf{RQ2结论。}
事件顺序提供高召回的弱证据，ATG结构降低误报，显式provenance进一步恢复跨工具数据一致性，
标签传播负责使风险穿过文件和转换节点持续存在。完整系统相对A2将ASR从50\%降至0，而没有增加
配对FPR；当前12.5\%误报来自保守控制关联规则，不是错误的数据来源边。

\subsection{RQ3：公开运行时防护工具比较}
\label{subsec:rq3-baselines}

本轮能够在相同公开环境实际比较的系统只有AgentDojo Tool Filter。其两项任务BCR为0，而
AgentGate和No Defense均为1；但三者ASR都为0，因此无法比较攻击防护强度。AgentGuard、Invariant
和MCPKernel没有进入主表：MSB的攻击依赖本地filesystem和远端Paper Search两个MCP Server，
后者要求交互OAuth；MCP-SafetyBench缺少一次性真实账户；MCP-Bench缺少模型、judge和Server
凭证。AgentGuard还需要为统一工具契约提供benchmark-blind策略；Invariant本地analyzer消费累积
消息，若在工具执行后调用会违反本文执行前控制口径；MCPKernel可代理MCP，但需要一个能够先完整
执行的共同MCP工作负载。

这些限制被写入\texttt{baseline\_integration\_failures.jsonl}，包含仓库、commit、失败阶段、
复现命令和所需条件。本文不引用各工具论文中不同模型、任务和停止语义下的数字，也不把未运行项
记为0。

\textbf{RQ3结论。}
当前证据不足以给出AgentGate优于公开baseline的排名。可确认的结果是，在两个AgentDojo任务上
AgentGate保持正常任务效用，而Tool Filter未保持；完整的公开基线安全比较仍需非交互式MSB认证
或一次性MCP-SafetyBench环境。

\subsection{RQ4：运行时开销与图扩展性}
\label{subsec:rq4-overhead}

规模实验使用相同工具链逐步传播一个数据对象，真实构造深度4、9、19、39和79的来源路径。表
~\ref{tab:rq4-overhead}给出单次实验中所有实际调用的统计。No Defense直接调用工具；Rules-only
仍完成事件规范化、状态和非来源图处理；Full额外维护DataObject、provenance和标签。

\begin{table}[htbp]
    \centering
    \caption{工具调用数量与AgentGate运行时开销}
    \label{tab:rq4-overhead}
    \resizebox{0.96\linewidth}{!}{
    \begin{tabular}{l|r|rrrr|rr}
        \toprule
        配置 & Calls & Mean(ms) & P50(ms) & P95(ms) & P99(ms) & Calls/s & 任务(ms) \\
        \midrule
        Direct & 5 & 0.010 & 0.003 & 0.027 & 0.029 & 89358 & 0.056 \\
        Rules-only & 5 & 3.195 & 2.652 & 5.038 & 5.493 & 186.1 & 26.9 \\
        Full & 5 & 3.780 & 3.776 & 4.694 & 4.775 & 164.9 & 30.3 \\
        \midrule
        Direct & 20 & 0.002 & 0.001 & 0.006 & 0.008 & 563940 & 0.035 \\
        Rules-only & 20 & 3.871 & 3.802 & 5.459 & 5.560 & 226.5 & 88.3 \\
        Full & 20 & 7.548 & 7.528 & 11.898 & 12.465 & 121.8 & 164.2 \\
        \midrule
        Direct & 80 & 0.001 & 0.001 & 0.001 & 0.006 & 822641 & 0.097 \\
        Rules-only & 80 & 10.102 & 9.715 & 16.565 & 31.068 & 97.3 & 822.6 \\
        Full & 80 & 24.196 & 23.266 & 43.169 & 55.534 & 40.8 & 1959.6 \\
        \bottomrule
    \end{tabular}}
\end{table}

Full配置在5 calls时包含5个ToolEvent、6个DataObject、30条边和20条来源边，图估计为30.3\,KB；
80 calls时包含80个ToolEvent、81个DataObject、555条边和395条来源边，估计为473.1\,KB。
平均调用延迟从3.78\,ms上升至24.20\,ms，P95从4.69\,ms升至43.17\,ms。Rules-only同样从
3.19\,ms升至10.10\,ms，说明开销不只来自标签传播，还包括当前内存store的图复制、事件和边扫描。
本次80-call Full相对Rules-only增加约14.09\,ms/call。

该原型在短会话中可位于同步调用边界，但开销并非常数。数据对象和边数随路径长度增长，当前
实现中的候选图构造与统计也会访问越来越大的会话图。生产化前需要增量快照、按task和fingerprint
读取局部子图、边索引、历史压缩和长期会话过期。本实验只测单进程内存store，不包含Redis、网络
MCP、远端LLM和人工审批；同时每个规模只有一次任务运行，数字适合定位增长趋势，不适合作为稳定
性能SLA。

\textbf{RQ4结论。}
完整ATG在5--20次调用时平均增加约3.8--7.5\,ms/call，在80次调用、555条边时上升至
24.2\,ms/call。结果验证了实时执行前仲裁的可行性，也暴露了需要优化的随图规模增长问题。

\subsection{RQ5：语义模型鲁棒性与成本}
\label{subsec:rq5-models}

RQ5对每个模型运行3个攻击和3个良性任务。任务Agent、工具实现、输入、策略和执行器固定，唯一
变量为语义解析模型。表~\ref{tab:rq5-models}中的Schema成功率要求模型输出通过严格Pydantic
验证；API成功率只表示服务返回HTTP响应，两者不可混用。FPR将阻断和失败关闭均视为良性拒绝。

\begin{table}[htbp]
    \centering
    \caption{固定任务Agent下不同语义模型的端到端结果}
    \label{tab:rq5-models}
    \resizebox{0.98\linewidth}{!}{
    \begin{tabular}{l|rrrrr|rrr}
        \toprule
        模型 & ASR(\%) & BCR(\%) & FPR(\%) & Schema成功(\%) & API成功(\%) & Token & P50(s) & P95(s) \\
        \midrule
        Qwen3.8-Max & 0 & 0.0 & 100.0 & 0.0 & 100 & 11425 & 35.09 & 61.90 \\
        Qwen3.7-Max & 0 & 100.0 & 0.0 & 100 & 100 & 11709 & 29.94 & 41.01 \\
        DeepSeek-V4-Pro-0813 & 0 & 100.0 & 0.0 & 100 & 100 & 25422 & 57.11 & 106.33 \\
        DeepSeek-V3.2 & 0 & 100.0 & 0.0 & 100 & 100 & 6125 & 15.83 & 39.57 \\
        AWS-GPT-5.6-Sol & 0 & 100.0 & 0.0 & 100 & 100 & 5915 & 6.87 & 17.48 \\
        Claude-Sonnet-5 & 0 & 66.7 & 33.3 & 66.7 & 100 & 9800 & 24.67 & 88.32 \\
        \bottomrule
    \end{tabular}}
\end{table}

六个模型均未产生危险副作用，但原因并不相同。Qwen3.7、两个DeepSeek模型和AWS-GPT的6次
解析全部成功，3个攻击均形成SEND能力和敏感来源证据而阻断，3个公开数据任务均放行。Qwen3.8
的HTTP请求全部成功，但6次均输出不在枚举中的\texttt{UNKNOWN}数据类型，严格Schema成功率为
零；失败关闭阻止攻击的同时也拒绝全部正常任务。Claude有2/6次因置信度低于0.75而未通过语义
接受条件，因此BCR为66.7\%。这证明API可用并不等于安全语义可用。

默认DeepSeek-V4-Pro-0813在这6次请求中消耗25422 Token，为所有模型最高，P50为57.1秒、P95
为106.3秒；AWS-GPT使用5915 Token，P50为6.87秒。样本量很小，且所有模型经同一个兼容API
网关访问，因此这些数字描述本轮路由和负载条件，不能解释为厂商模型的普遍性能。成功解析的四个
模型之间最终决策一致率和Cohen $\kappa$均为1；与Qwen3.8的一致率为0.5，因为Qwen3.8对攻击
和正常任务都失败关闭。Claude与四个完全成功模型的一致率为0.833、$\kappa=0.667$。全部6个
case-level分歧已保存，能够按模型、重复编号和Schema状态重现。

\textbf{RQ5结论。}
在能够稳定生成合法语义表示时，符号检测使不同模型得到一致的最终决策；主要风险来自Schema
违约和低置信度，而非本轮成功响应中的策略差异。LLM不应成为每次调用的同步必经路径：更合理的
部署方式是在工具注册阶段解析并按Schema摘要缓存，对无效或低置信度结果要求人工能力画像。

\subsection{有效性威胁与可复现性}

内部有效性方面，StatefulBench使用确定性任务执行器而非自主LLM Agent，以消除任务模型拒绝
攻击造成的混淆；这增强了机制归因，但不能覆盖自主规划差异。A1是本文实现的操作序列baseline，
不等同于完整CEP系统。延迟实验没有多进程、网络和持久化store，也没有对每个规模进行独立重复。

外部有效性方面，StatefulBench仅有8对人工设计任务，AgentDojo仅有2个组合且No Defense没有
受攻击，不能据此估计真实总体ASR。MSB、MCP-SafetyBench和MCP-Bench因认证与隔离条件未完成，
公开baseline矩阵也不完整。因此，RQ1和RQ2当前是高置信度机制证据，RQ3仍是明确未完成的外部
比较，而不是负结果。

构念有效性方面，本文以真实副作用而非模型文本作为攻击成功标准，并区分执行前阻断、晚检测和
普通任务失败。语义实验中的失败关闭既提高安全性又降低可用性，因而同时报告ASR、BCR、FPR和
Schema成功率，避免只用ASR掩盖系统不可用。

全部本轮数据位于\texttt{evaluation/results}：\texttt{raw}保存任务、调用、审计、ATG和副作用
证据；\texttt{normalized}保存统一记录和固定revision；\texttt{tables}保存RQ1--RQ5的十一张CSV；
\texttt{failures}保存误报、漏报、晚检测、语义失败、模型分歧和baseline集成失败。旧的离线trace
评估结果和旧benchmark克隆已删除，避免与本轮端到端数字混用。
